% STATUT : À RÉDIGER (à faire en dernier, une fois les résultats figés).
\chapter{Introduction}
\label{ch:introduction}

\textit{[Chapitre à rédiger.]}

% PLAN OBLIGATOIRE :
%  1. La question : DORA impose la conformité de cinq piliers, mais aucun
%     cadre ne traduit la non-conformité en capital. La Formule Standard
%     n'exprime rien de l'effet DORA (0 %, cf. chapitre benchmarks).
%  2. L'intuition : une défaillance ne reste pas dans son pilier, elle se
%     propage, et elle se propage DANS UN SENS. Exemple concret.
%  3. La contribution, en trois points :
%       (i)  un mécanisme de dépendance dirigée et à l'ordre entre
%            domaines de contrôle, avec sa non-transitivité ;
%       (ii) une frontière d'identifiabilité démontrée, muée en
%            spécification de reporting DORA ;
%       (iii) une traduction en capital, présentée en distribution
%            et en bornes, jamais en point.
%  4. Ce que le mémoire NE prétend PAS : W n'est pas calibré, le niveau
%     absolu du SCR est illustratif, c'est l'effet relatif qui est vendu.
%  5. Plan du document.
